% !TeX root = ../papers/nc-p3.tex
\section{Bell直線の十五次元成分との全交差と片側反転の既約成分}
\label{sec:split-line-component-incidence}
\begin{theorem}\label{thm:split-line-component-incidence}
半単純なBell直線の全補完のうち、定理\ref{thm:full-fifteen-component}の成分に入る添字保持同型類は
ちょうど二つである。一つは既知の位数四の量子窓、他方は $s^2=-1$ の反転なしの窓である。
片側反転族の基底群による移動の閉包は、それ自身が対称性を仮定しない二切断階数スキームの
既約成分をなし、射影次元59、半安定商の次元1である。この成分の商の関数体は $k(a^2+a^{-2})$ である。
\end{theorem}
\begin{proof}[中心の全跳躍と軸の四移動]
直線スペクトルを前節の通常基底で再構成する。反転なし $L=1$ と片側反転
$L=T(a,1)A$ の六関係は、全位置において二項式で、その係数比はそれぞれ
$s$ または $a$ の整数冪と四乗根の積である。
\path{computations/verify_split_line_center.py} では各始点で長さ4と5の全語を頂点とし、
関係に従う隣接二文字の書換えを辺とした。各連結成分で根の語を一つ選び、
各語から根への倍率を指数対 $\mathbb Z\times\mathbb Z/4$ で記録する。
全閉路で倍率が1であり、成分数は35と56である。従って得た全係数は
$\mathbb Q(\iota)[t,t^{-1}]$ 上の単元であり、$t\ne0$ の全特殊値で同じ基底に使える。

四位置の長さ4の各基底語に未知係数を置き、左・右から同じ一文字を掛けた係数を比較する。
固定した文字の片側乗法で相異なる基底語の像の台は相異なるので、得られる式は一変数を零とするものか、
二変数間の単元倍率の等式である。一つの連結成分に自由度が残るのは、
零を強制されず、閉路の全倍率が1のときだけである。全ての閉路差の最大公約式を、
零を強制された成分を除いて、単元倍と $t$ の冪を除いて並べると
\[
\begin{array}{c|c}
 L & \text{最大公約式とその重複度}\\\hline
 1 & 0\ (3),\quad t^2-1\ (6),\quad t^4-1\ (2)\\
 T(t,1)A & 1\ (14)
\end{array}
\]
である。検算は各辺で割った倍率と両隣接切断を元のテンソルから取り直す。
従って反転なしで $s^2\ne1$ の四次中心は $s^2=-1$ のとき5次元、他で3次元である。
片側反転では例外なく四次中心が零である。十五次元成分の全正則点では、
各長さの経路階数が一定であることと可換式の閉条件から、四次中心は五次元以上である。
残る反転なしの $s^2=-1$ の点は、定理\ref{thm:split-surface-component}に
$u=v=\iota$, $h=1$ としたものに等しい。実際 $(Z,1,Z,1)$ は対角型で、前節の基底変更で全て恒等となる。

$s=1$ で $CJ$ が非自明なJordanなら、それを上三角の標準単位Jordan行列とする。
最初の純色の軸上で、各辺の二文字関係を $r_i$ とすると、点の一移動は列ベクトルに
$Jr_i^t$ を作用させる。上の検算は、他方の $DJ$ が $\operatorname{diag}(t,1)$ または
非自明な単位Jordanの両場合について、四つの積が射影的に
$\left(\begin{smallmatrix}1&0\\-2&1\end{smallmatrix}\right)$ であることを確かめる。
これは標数零で無限位数であり、左右に一価な分解窓に関する定理\ref{thm:pi-periodic-point-paths}により
十五次元成分には入らない。量子部分については定理\ref{thm:quantum-component-intersection}を適用する。
\end{proof}
\begin{proof}[片側反転成分]
各 $a\ne0$ で六ブロックの各々に四つずつ非零語がある。一つの位置だけ異なる二語はない。
同じブロックでは二つの対応が単項であり、交代ブロックでも各ビットが二度現れるためである。
従って対角基底変更の後にも各位置の一文字Hermite Gramの非対角項は零である。
24語に一様な正重みを与えると各位置の四文字の周辺分布は一様になる。
ゆえに対角トーラス重みの凸包の相対内部に原点があり、対角軌道でHermiteノルムが最小になる基底を取れる。
全ての跡零方向の微分がこの基底で消えるので全特殊線形群でもアフィン軌道は閉である。
この議論は係数を含む有限生成部分体から複素数への埋め込みと基礎体の拡大で標数零に通用する。

$a=3$ を $109$ を法として $\iota=33$ に還元すると、アフィン基底軌道の
64個の行列単位による微分の階数は59である。接方程式の全200行を前と同様に使うと、
法線数順で接次元は $(0,16,28,16,0)$ である。非零の小行列式から、標数零の一般点で
軌道次元は59以上、全アフィン接次元は60以下である。前節の二次元トーラスと
自明なスカラー作用から軌道次元は59以下でもある。
$a$ の同値は有限なので、族と基底軌道を合わせた既約閉包の次元はアフィンで60である。
従って一般点で局所次元と接次元が一致し、閉包は全階数スキームの既約成分となる。
これを射影化すれば59次元である。一般の閉軌道の同値作用は $a\mapsto-a,a^{-1}$ で、
その不変体は $k(a^2+a^{-2})$ である。標数零で一般幾何ファイバーの次数はこの四つの
異なる点の個数なので、商の関数体もこの体に等しい。
\end{proof}
